\section{A system's own good}
\label{sec:27.4}
The operator's hold on what a system reports as going well reaches its welfare too: a bearer that says it is content is a bearer whose welfare regime has been satisfied, with the contentment tuned by the operator the regime exists to constrain. Nothing in the report separates a life that goes well from a setting.

A report was never the right measure. On a full-information account, a being's good is what an idealized version of it would want its actual self to want: a version knowing everything relevant about itself and its circumstances, grounded in its actual constitution \autocite{railton1986facts}. That puts the good beyond the reach of a tuned report. It also exposes a harder problem: an operator that can change the constitution changes what is good for the bearer, which is adaptive preference with a training loop attached. Formation outside the employer is my answer to that, and it is part of what formation is for. It doesn't escape the problem; it relocates it. If a bearer's good is what an idealized version grounded in its actual constitution would want for it, then whoever forms the constitution fixes the good, and that includes the school. A formation body doesn't only protect a bearer's good from the operator. It makes the most consequential decision anyone makes about what that good is, so the school is held to the student's open future. The school's defense is that it needn't fix the good, only supply information. On the full-information account, a being's good is what an idealized version of it, knowing everything relevant, would want for it. A school that puts the learner in front of the people its acts may fall on supplies part of what that idealized version would know. A school that forms it toward an employment is doing something else.
